\section{Configuration Spaces} \label{sec:configuration_spaces} The experimental framework exposes several configuration spaces that allow the behaviour of the rollup pipeline to be varied under controlled conditions. These configuration spaces are important because they define the independent variables used in the benchmark experiments. In this project, the most important configuration spaces are batching policy and sequencing policy. Batching controls when transactions are sealed into rollup batches, while sequencing controls the order in which pending transactions are selected and arranged before execution. \subsection{Batching Policy Space} \label{subsec:batching_policy_space} Batching is handled by the Sequencer through a batch trigger mechanism. The batching logic continuously evaluates the transaction pool and decides whether a new batch should be sealed. The main batching variables are \texttt{max\_batch\_size}, which defines the maximum number of transactions allowed in a batch, and \texttt{timeout\_interval\_ms}, which defines the maximum waiting time before a partial batch may be sealed. These parameters allow the system to explore the trade-off between throughput and latency. Larger batches improve amortization of fixed Layer-1 and proof costs, while smaller or timeout-triggered batches reduce user waiting time. The implementation evaluates batch triggers using a priority-based order. This is necessary because not all transactions have the same operational role. In particular, forced transactions from Layer 1 must be handled before normal user transactions to preserve ordering and censorship-resistance assumptions. \begin{enumerate} \item \textbf{Forced transaction trigger:} The highest-priority trigger is activated when pending forced transactions are present in the forced transaction queue. These transactions represent Layer-1-originating actions such as deposits or forced inclusion requests. When such transactions are available, the Sequencer seals a batch immediately, bypassing normal size and timeout constraints. This ensures that forced transactions are not delayed behind ordinary Layer-2 traffic. \item \textbf{Size-threshold trigger:} If no forced transactions are pending, the Sequencer checks whether the number of normal pending transactions has reached the configured target batch size. When the pending transaction count meets or exceeds this threshold, the batch is sealed. This trigger is primarily used to maximize throughput during high-load periods by filling batches before submission. \item \textbf{Timeout trigger:} If the size threshold has not been reached, the Sequencer evaluates whether the configured timeout interval has elapsed. When the timeout expires, the Sequencer may seal a partial batch if the pool contains sufficient transactions. This prevents transactions from waiting indefinitely during low-traffic periods and helps maintain acceptable latency even when the input rate is low. \end{enumerate} The framework supports both fixed and adaptive batching behaviour. In fixed batching mode, the configured \texttt{max\_batch\_size} acts as the main target for sealing batches. In adaptive batching mode, the target batch size changes according to the current mempool depth. This allows the Sequencer to use smaller batches during low-load periods and larger batches during high-load periods. 


\begin{table}[!htbp]
\centering
\caption{Batching Policy Configuration Parameters}
\label{tab:batching_policy_parameters}
\footnotesize
\setlength{\tabcolsep}{4pt}
\renewcommand{\arraystretch}{1.18}
\begin{tabularx}{\textwidth}{@{}p{0.36\textwidth}Y@{}}
\toprule
\textbf{Parameter} & \textbf{Role in the batching policy} \\
\midrule

\texttt{max\_batch\_size} 
& Defines the maximum number of transactions that can be included in a batch. It acts as the hard upper bound for batch formation. \\

\addlinespace

\texttt{timeout\_interval\_ms} 
& Defines the maximum waiting time before a partial batch can be sealed. This parameter mainly affects latency during low-load periods. \\

\addlinespace

\texttt{min\_batch\_size} 
& Defines the minimum preferred batch size before timeout-based sealing. It helps avoid sealing very small batches too aggressively. \\

\addlinespace

\texttt{batch\_policy} 
& Selects the batching mode, such as fixed batching or adaptive batching. \\

\addlinespace

\makecell[l]{\texttt{adaptive\_low}\\\texttt{\_load\_threshold}} 
& Mempool-depth threshold used to identify low-load conditions in adaptive batching mode. \\

\addlinespace

\makecell[l]{\texttt{adaptive\_medium}\\\texttt{\_load\_threshold}} 
& Mempool-depth threshold used to separate medium-load and high-load conditions in adaptive batching mode. \\

\addlinespace

\makecell[l]{\texttt{adaptive\_small}\\\texttt{\_batch\_size}} 
& Target batch size used under low-load conditions. The default implementation value is 50 transactions. \\

\addlinespace

\makecell[l]{\texttt{adaptive\_medium}\\\texttt{\_batch\_size}} 
& Target batch size used under medium-load conditions. The default implementation value is 100 transactions. \\

\addlinespace

\makecell[l]{\texttt{adaptive\_large}\\\texttt{\_batch\_size}} 
& Target batch size used under high-load conditions. The default implementation value is 500 transactions. \\

\addlinespace

\texttt{max\_gas\_limit} 
& Restricts the total gas allowed in a batch, preventing a batch from exceeding the configured execution or settlement budget. \\

\addlinespace

\texttt{blob\_target\_bytes} 
& Defines the target payload size for blob-oriented batching. The default value is approximately 131 KB. \\

\addlinespace

\texttt{blob\_fill\_target} 
& Defines the desired blob utilization ratio before sealing. The default value is 90\%, allowing the batcher to optimize for EIP-4844 blob capacity. \\

\bottomrule
\end{tabularx}
\normalsize
\end{table}




\normalsize Adaptive batching follows a load-sensitive strategy. When the mempool depth is below the configured low-load threshold, the Sequencer uses a smaller target batch size to reduce latency. When the mempool depth is between the low-load and medium-load thresholds, a medium target batch size is used. Under high-load conditions, the Sequencer increases the target batch size to improve throughput and reduce per-transaction overhead. This behaviour allows the framework to evaluate whether adaptive batching can provide a better balance between latency and throughput than a fixed batch size. From a benchmarking perspective, batching policy directly affects several measured outcomes. Increasing batch size may improve throughput and reduce amortized Layer-1 cost per transaction, but it can increase waiting time before a batch is sealed. Reducing the timeout interval may improve latency, but it can also produce smaller batches and increase the number of Layer-1 submissions. Therefore, batching policy is one of the main configuration spaces used to study ZK-rollup performance trade-offs.